\documentclass{article}
\usepackage[UTF8]{ctex}
\usepackage{amsmath,amssymb,amsfonts}
\usepackage{enumitem}
\usepackage{geometry}
\geometry{a4paper, margin=1in}
% === 错题本排版优化（不参与题目标识） ===
\setlength{\parindent}{0pt}
\setlength{\parskip}{0.35em plus 0.10em minus 0.08em}
\setlength{\emergencystretch}{3em}
\setlist{itemsep=0.25em, topsep=0.40em, parsep=0pt}
\allowdisplaybreaks[2]
% === 错题本排版优化结束 ===
\begin{document}
\section*{福州大学2025年高等代数真题}
\section*{109．福州大学2025年高等代数试卷}
一．填空题．\\

\subsection{题目1}
1．设 \(\alpha_{1}=(0,1,2)^{\mathrm{T}}, \alpha_{2}=(-1,3,2)^{\mathrm{T}}, \alpha_{3}=(1, \lambda, 3)^{\mathrm{T}}\) ，当且仅当 \(\lambda\) 满足 \(\_\_\_\_\)时，任意3维列向量都可由 \(\alpha_{1}, \alpha_{2}, \alpha_{3}\) 线性表出。

\subsection{题目2}
2．设 \(A\) 是秩为 \(r\) 的 \(n\) 阶方阵，且 \(A^{2}=-A, E\) 为 \(n\) 阶单位矩阵，则 \(|2 E+A|=\) \(\_\_\_\_\)．

\subsection{题目3}
3．设 \(A, B\) 分别为 \(n \times m\) 和 \(m \times n\) 矩阵，若 \(r(A B)=n\) ，则 \(r(B A)=\) \(\_\_\_\_\)．

\subsection{题目4}
4．设 \(A=\alpha \beta^{\mathrm{T}}\) ，其中 \(\alpha=(1,3,4)^{\mathrm{T}}, \beta=(2,2,1)^{\mathrm{T}}\) ，则 \(A\) 的 Jordan 标准形为 \(\_\_\_\_\)．

\subsection{题目5}
5．设 \(A\) 为 3 阶正交矩阵，且 \(a_{22}=-1\) ，则 \(A X=(0,1,0)^{\mathrm{T}}\) 的解为 \(\_\_\_\_\)．

二．简答题。\\

\subsection{题目6}
6．\(A\) 是 \(m \times n\) 矩阵，证明：存在 \(n \times m\) 矩阵 \(B\) ，使得 \(A B A=A\) 且 \(B A B=B\) 。

\subsection{题目7}
7．设 \(A=\left(\begin{array}{llllll}1 & 0 & 0 & 0 & 0 & 0 \\ 0 & 1 & 0 & 0 & 0 & 0 \\ 0 & 1 & 1 & 0 & 0 & 0 \\ 0 & 0 & 0 & 2 & 0 & 0 \\ 0 & 0 & 0 & 1 & 2 & 0 \\ 0 & 0 & 0 & 0 & 1 & 2\end{array}\right)\) ，求 \(A\) 的不变因子组和初等因子组。

\subsection{题目8}
8．设

\[
F[x]_{n}=\left\{a_{n} x^{n}+\cdots+a_{1} x+a_{0} \mid a_{n}, \cdots, a_{1}, a_{0} \in F\right\} .
\]

求线性变换 \(\mathcal{D}: F[x]_{n} \rightarrow F[x]_{n}, \mathcal{D}(f(x))=f^{\prime}(x)\) 的特征多项式和所有特征值．

\subsection{题目9}
9．设 \(\alpha\) 为 \(n\) 维实单位列向量，\(E\) 为 \(n\) 阶单位矩阵，求 \(E-\alpha \alpha^{\mathrm{T}}\) 的秩并说明理由．

\subsection{题目10}
10．\(A=\left(\begin{array}{ccc}1 & 4 & -2 \\ 0 & -1 & 0 \\ 1 & 2 & -2\end{array}\right)\) ，求 \(A^{2025}\) ．

三．问答题．\\

\subsection{题目11}
11．用正交线性替换将二次型 \(f\left(x_{1}, x_{2}, x_{3}\right)=2 x_{1}^{2}+2 x_{2}^{2}+2 x_{3}^{2}-2 x_{1} x_{2}-2 x_{1} x_{3}-2 x_{2} x_{3}\) 化为标准形，写出所做的正交线性变换和得到的标准形。\\

\subsection{题目12}
12．设 3 阶实对称矩阵 \(A\) 的各行元素之和均为 \(4, ~ r(A-E)=1\) ，其中 \(E\) 为3阶单位矩阵。求矩阵 \(A\)．\\

\subsection{题目13}
13．设 \(n\) 阶方阵 \(A, B\) 满足 \(A B=B A\) ，证明：\(\left|\begin{array}{ll}A & B \\ B & A\end{array}\right|=\left|A^{2}-B^{2}\right|\) 。\\

\subsection{题目14}
14．设 \(A, B, A B-E\) 都是 \(n\) 阶可逆阵，\(E\) 为 \(n\) 阶单位矩阵。

\begin{itemize}
  \item[（1）] 证明：\(A-B^{-1}\) 可逆，并求其逆．
  \item[（2）] 证明：\(\left(A-B^{-1}\right)^{-1}-A^{-1}\) 可逆，并求其逆．
\end{itemize}

\subsection{题目15}
15．设 \(A\) 为 \(n\) 阶实反对称矩阵，\(\lambda\) 是 \(A\) 的特征值，证明：\(\lambda\) 实部为 0 。\\

\subsection{题目16}
16．设 \(\varphi\) 为线性空间 \(V\) 上的线性变换，若 \(\varphi\) 的特征多项式为 \(f(\lambda)=\lambda^{k}(\lambda-1)^{n-k}\) ，其中 \(0<k<n\) 。证明：\(V=V_{1} \oplus V_{2}\) ，其中 \(V_{1}=\operatorname{Ker} \varphi^{k}, V_{2}=\operatorname{Ker}\left(\varphi-\operatorname{id}_{V}\right)^{n-k}\) 。\\

\subsection{题目17}
17．设 \(\varphi, \psi\) 为线性空间 \(V\) 上的线性变换，且 \(\varphi^{2}=\varphi\) 。证明：

\begin{itemize}
  \item[（1）] \(\operatorname{Ker} \varphi=\{\alpha-\varphi(\alpha) \mid \alpha \in V\}\) ．
  \item[（2）] \(\operatorname{Ker} \varphi, \operatorname{Im} \varphi\) 均是 \(\psi\)－不变子空间的充分必要条件是 \(\psi \varphi=\varphi \psi\) 。
\end{itemize}

\subsection{题目18}
18．设 \(A\) 为 \(m \times n\) 实矩阵。证明：

\begin{itemize}
  \item[（1）] 对任意 \(m\) 维实向量 \(\beta\) ，线性方程组 \(A^{\mathrm{T}} A X=A^{\mathrm{T}} \beta\) 必有解。
\end{itemize}

（2）若 \(X_{0}\) 是线性方程组 \(A^{\mathrm{T}} A X=A^{\mathrm{T}} \beta\) 的解，则在欧氏空间 \(\mathbb{R}^{n}\) 中，对任意 \(\alpha \in \mathbb{R}^{n}\) ，均有\(d\left(\beta, A X_{0}\right) \leq d(\beta, A \alpha)\) ，其中 \(d(X, Y)\) 表示向量 \(X, Y\) 的距离．
\end{document}
